设 \(X\) 为 \(k\) 上的形式簇。若 \(x\in X\) 投影到 \(s\in\operatorname{Spf}(k)\)，则剩余域 \(\kappa(x)\) 为 \(\kappa(s)\) 的有限扩张。以 \(\kappa_e(x)\) 表示 \(\kappa(s)\) 在 \(\kappa(x)\) 中的可分闭包。\footnote{\label{sga3:VIIB:nde:73}下文中，编者利用前面的增补，详细给出函子 \(X\mapsto X_e\) 的构造，并证明它与有限乘积交换；\SGARefLink{sga3:VIIB:2.5.1:corollary}{2.5.1} 将用到这一点。}

\medskip\noindent\textbf{命题 1.6.F。}\phantomsection \label{sga3:VIIB:1.6.F}{\itshape \SGARefLink{sga3:vii:paragraph:5-1:equation:i}{(i)} 包含函子 \[
  \operatorname{Vaf}^{\mathrm{\acute{e}t}}/k
  \longrightarrow
  \operatorname{Vaf}/k
\] 有左伴随 \(X\mapsto X_e\)。形式簇 \(X_e\) 与 \(X\) 有相同的点。对每个投影到 \(s\in\operatorname{Spf}(k)\) 的 \(x\in X\)，选取 \(\kappa_e(x)/\kappa(s)\) 的本原元 \(\xi\)，设其次数为 \(n\)，任取 \(\xi\) 的提升 \(u\in\mathcal O_{X,x}\)，并设 \(F\in k[T]\) 为次数为 \(n\)、零化 \(u\) 的首一多项式。然后置 \[
  \mathcal O_{X_e,x}=k[T]/(F).
\]。对每个 \(Y\in\operatorname{Ob}(\operatorname{Vaf}^{\mathrm{\acute{e}t}}/k)\)，下图中的典范映射均为双射：\begin{center}\phantomsection \label{sga3:VIIB:1.6.F:hom-square} % SGA3-VII-DIAG-095
% Authority: Exp7B-13oct24.pdf, local p. 42.
% Crop: figures/Exp7B_localp042_etale_reflection_hom_square.png
% Locator: tex/components/34_expose_VIIB_section16F_etale_reflection_en.tex, Loop 1 line 45.
\begin{tikzcd}[row sep=2.8em,column sep=3.5em]
\operatorname{Hom}_{\operatorname{Vaf}/k}(X,Y)
  \arrow[r,"{\sim}"]
  & \operatorname{Hom}_{\operatorname{Vaf}/\kappa(k)}
      (X_\kappa,Y_\kappa)
      \arrow[d,"\wr"] \\
\operatorname{Hom}_{\operatorname{Vaf}/k}(X_e,Y)
  \arrow[r,"{\sim}"]
  & \operatorname{Hom}_{\operatorname{Vaf}/\kappa(k)}
      ((X_e)_\kappa,Y_\kappa)\, ,
\end{tikzcd}
\end{center}竖箭头由包含映射 \[
  \kappa_e(x)\lhook\joinrel\longrightarrow\kappa(x),
  \qquad x\in X.
\] 诱导。特别地，这定义了态射 \[
  p:X\longrightarrow X_e,
\]，且对每个形式 \(k\)-簇态射 \(\phi:X\to Y\)，若 \(Y\) 在 \(k\) 上 平展，则该态射唯一地通过 \(p\) 分解。

\SGARefLink{sga3:vii:paragraph:5-1:equation:ii}{(ii)} 函子 \(X\mapsto X_e\) 与有限乘积交换。}

\medskip\noindent\textit{证明。}断言 \textup{\SGARefLink{sga3:vii:paragraph:5-1:equation:i}{(i)}} 由 \SGARefLink{sga3:VIIB:1.6.E}{1.6.E}~\textup{(b)} 和 \SGARefLink{sga3:VIIB:1.6.D}{1.6.D} 得出。现证明 \textup{\SGARefLink{sga3:vii:paragraph:5-1:equation:ii}{(ii)}}。由 \SGARefLink{sga3:VIIB:1.6.E}{1.6.E}~\textup{(b)} 中的范畴等价，可以假设 \(k\) 是域。此时，若 \(X\) 是半局部形式 \(k\)-簇，即其仿射代数 \(A\) 为半局部环，则 \(X_e\) 的仿射代数就是 \(A\) 中在 \(k\) 上 平展 的最大子代数；将它记为 \(A_e\)。

因此只需证明：若 \(K\) 和 \(L\) 为 \(k\) 的两个有限扩张，则 \[
  K_e\otimes_kL_e
  =
  (K\otimes_kL)_e.
\]。设 \(p\) 为 \(k\) 的特征指数，即当 \(\operatorname{char}(k)=0\) 时 \(p=1\)，否则 \(p=\operatorname{char}(k)\)。对每个 \(x\in K\otimes_kL\)，存在整数 \(n\geq1\) 使得 \[
  x^{p^n}\in K_e\otimes_kL_e.
\]。于是 \(K\otimes_kL\) 在 \(K_e\otimes_kL_e\) 上是纯不可分的（radicial），且其最大 平展 \(k\)-子代数恰为 \(K_e\otimes_kL_e\)。故 \[
  K_e\otimes_kL_e=(K\otimes_kL)_e.
\]。

沿用此记号，注意 \(\mathcal O_{X_e,x}\) 未必是 \(\mathcal O_{X,x}\) 的子代数。由下一命题，若 \(X\) 在 \(k\) 上拓扑平坦，则它是子代数。
